% 附录一 雍亲王致年羹尧书
% 原书页码：336；标题
\BookUnnumbered{附录一\quad 雍亲王致年羹尧书}
\par
% 原书页码：336；段落
自《文献丛编》（北京，故宫博物院，1930），卷1。这封信的原%
件图片在全国报刊索引和读秀的数据库上都能调阅。%
\par
% 原书页码：336；段落
{\bookkai 知汝以儇佻恶少，屡逢侥幸，君臣大义，素所面墙。国朝祖宗制}%
{\bookkai 度，各王门旗属，主仆称呼，永垂久远，俱有深意。尔狂昧无知，具}%
{\bookkai 启称职，出自何典？屡谕尔父，尔犹抗违不悛，不徒腹诽，而竟公然}%
{\bookkai 饰词诡拒，无父无君，莫此为甚！}%
\par
% 原书页码：336；段落
{\bookkai 况妃母千秋大庆，阿哥完婚之喜，而汝从无一字前来称贺，六七}%
{\bookkai 个月无一请安启字，视本门之主已同陌路人矣。且汝所称，捐资助饷，}%
{\bookkai 家无余财，更属无谓之甚。况我从未问及汝家囊橐，汝何得以鄙亵之}%
{\bookkai 心测我，而肆进其矫廉之词？}%
\par
% 原书页码：336–337；段落（跨页续段）
{\bookkai 况汝在蜀，骄横不法，狂悖无忌，皇上将来不无洞鉴。而尚敢谓}%
{\bookkai 今日之不负皇上，即异日之不负我者，是何言欤。以无法无天之谈而}%
{\bookkai 诱余以不安分之举也，岂封疆大臣之所当言者？“异日”两字足可以}%
{\bookkai 诛羹尧全家。且汝于孟光祖馈遗授受，不但众所共知，而且出自汝家}%
{\bookkai 人之亲口以告我者，尚敢矇眬皇上，得以漏网？即此一事，是即汝之}%
{\bookkai 现在所以负皇上而将来之所以必负我者也！}%
\par
% 原书页码：337；段落
{\bookkai 至于我之培植下人，即并其家人父子亦无不委曲作养成全，在汝}%
{\bookkai 固已无人心，谅必非无耳无目者。于此不思所以报称，而反公然跋扈，}%
{\bookkai 尔所蓄何心，诚何所挟持而竟敢于如此耶？即此无状，是即汝之现在}%
{\bookkai 所以负我，即异日必负皇上者也！}%
\par
% 原书页码：337；段落
{\bookkai 况在朝廷称君臣，在本门称主仆，故自亲王、郡王、贝勒、贝子}%
{\bookkai 以至公等，莫不皆称主子、奴才，此通行常例也。且汝父称奴才，汝}%
{\bookkai 兄称奴才，汝父岂非封疆大臣乎？而汝独不然者，是汝非汝兄之弟，}%
{\bookkai 亦非汝父之子矣？又何必称我为主？既称为主，又何不可自称奴才}%
{\bookkai 耶？汝父兄所为，不是，汝当劝；约而同之，则犹可也。不遵父训，}%
{\bookkai 抗拒本主，无父无君，万分可恶。若汝或另有所见，或别有委曲，汝}%
{\bookkai 不妨具折启奏，申明汝之大典。我亦将汝必不肯称奴才之故，以至妃}%
{\bookkai 母大庆、阿哥喜事，并于我处终年无一字请安，以及孟光祖之事，与}%
{\bookkai 汝所具“异日”之启，好好存留在此，一一奏明，谅皇上自有定夺也。}%
\par
% 原书页码：337；段落
{\bookkai 再，汝父年老，汝子自当代汝奉养。汝毫不为意，七八个尽留任}%
{\bookkai 所，岂人心之能恶也。只待汝子娶亲方令来京，信乎？求忠臣于孝子}%
{\bookkai 也。而又使及于我，所具启仪，苟简无礼，言词皆谬，皆汝之不肖下}%
{\bookkai 属，无可奈何之所以应塞汝者，而概施之于我，是岂主子奴才之礼乎？}%
{\bookkai 凡此者皆汝之不学无术，只知逞一时刚愎之私，而自贻乃父之戚耳。}%
\par
% 原书页码：337–338；段落（跨页续段）
{\bookkai 自今以后，凡汝子十岁以上者，俱着来京侍奉汝父，即汝昔年临}%
{\bookkai 行时向我讨去读书之弟侄，亦必着令作速来京，毋留在外，法成汝无}%
{\bookkai 父无君之行也！}%
\par
% 原书页码：338；段落
{\bookkai 观汝今日藐视本门主子之意，他日为谋反叛逆之举，皆不可定。}%
{\bookkai 汝父见汝此启，当余之面痛哭气恨倒地，言汝疯狂乱为。汝如此所为，}%
{\bookkai 而犹敢以伪孝欺人，蚬言父子天性，何其丧心病狂，一至于此？}%
\par
% 原书页码：338；段落
{\bookkai 况汝父在京，我之待他恩典甚重，谅汝无父之人亦未必深悉其}%
{\bookkai 委曲也。然圣主以孝治天下，而于我惜老之夙心有所不忍，故不惜}%
{\bookkai 如此申斥，警汝愚蒙。汝诚能于此爽然自失，真实悔悟，则诚汝之}%
{\bookkai 福也！其犹执迷不悛，则真所谓噬脐莫及者矣！汝其图之！}%
